\scanpage{91}\mspage{82}
\[
3^\circ)\quad X\,|\circ|\,Y\Longleftrightarrow
 \left\{\begin{aligned}
 X=P_x,Y=P_y:\quad &[x\ne y\ \text{i.e.}]\quad x<y\ \text{ou}\ y<x\\
 X=P_x,Y=S_{a,b}\ (\text{ou l'inverse}):\quad &x\leq a\ \text{ou}\ b\deleted{\unsure{<|>}}\leq x\\
 X=S_{a,b},Y=S_{a',b'}\quad &a<b\leq a'<b'\ \text{ou}\ a'<b'\leq a<b
 \end{aligned}\right.
\]

On a exprimé simplement $|X|^\circ\cap|Y|^\circ\deleted{\unsure{\subset|=}}=\varnothing$, en utilisant la divisibilité.

Dans le contexte ensembliste actuel, on n'a plus d'axiomes à vérifier, pas non plus \deleted{\unsure{des}}d'axiomes du support, grâce à l'existence de la fonction support externe fidèle $X\mapsto|X|$ (comme \deleted{\unsure{chaque}}\inserted{pour} \deleted{\unsure{fois}} tout magma ensembliste, ou simplement ponctuel, voire, seulement fidèle -- en prenant cette fois la fonction support fidèle $X\mathrel{\tikz[baseline=-.45ex,x=1em,y=1em]{\draw (0,0)--(1.15,0) (0,-.10)--(0,.11) (1.15,-.10)--(1.15,.11);}}\unsure{\mathrm{Omb}\,X|\mathrm{Ouvs}\,X|\mathrm{Oens}\,X}$).

B) Supposons maintenant la relation d'ordre quelconque. On prendra encore
\[
\left\{\begin{aligned}
 &M\subset\unsure{\mathcal P|\mathscr P}(\mathcal P(L))\\
 &M=M_0\amalg M_1\\
 &\qquad M_0=\{P_x\mid x\in L\},\qquad P_x=\{\{x\}\}\\
 &\qquad M_1=\{S_{a,b}\mid a,b\in L,\ a<b\}\qquad
 S_{a,b}=\{[a,b],\{a\},\{b\}\},
\end{aligned}\right.
\]
\margin{il \deleted{\unsure{est}} peut-être plus simple de prendre $M_0=L$, $M_1=\{(a,b)\in L\times L\mid a<b\}$}
mais on notera qu'on ne sait pas si $S_{ab}$ est une figure (élémentaire), faute de savoir si
\[
 |S_{a,b}|^\circ=[a,b]\setminus\{a,b\}=]a,b[
\]
\unsure{tjrs|tjous} non vide. On \uline{définira} alors les relations $\leq,\ \overset{\circ}{\ll},\ |\bullet|$ par les formules explicitées plus haut, en faisant attention cette fois que
\margin{\uline{NB} Avec la présentation adoptée, on trouve des ``quasi-figures'' élémentaires \unsure{dans|de} $L$ (on peut avoir $|X|^\circ=\varnothing$) et $\leq,\overset{\circ}{\ll}$ est induit alors, au sens évident, par les notions ensemblistes. Mais \uuline{pas} $|\circ|$, sauf dans le cas tot. ordonné.}
\[
\begin{aligned}
 S_{a,b}\overset{\circ}{\ll}S_{a',b'}\Longleftrightarrow a'\leq a<b\leq b'
 &\ \Longleftrightarrow |S_{a,b}|^{\deleted{\unsure{\circ}}}\overset{\inserted{\unsure{\nRightarrow|\Rightarrow}}}{\subset}|S_{a',b'}|^{\deleted{\unsure{\circ}}}\\
 &\quad\deleted{\text{\unsure{mais l'implication}}}\\
 &\quad\deleted{\text{\unsure{pas nécessairement valable}}}\\
 &\ \Longrightarrow |S_{a,b}|^\circ\subset|S_{a',b'}|^\circ
\end{aligned}
\]
\margin{mais l'implication inverse pas néc. valable}

\[
\deleted{P_x\,|\circ|\,P_y\Longleftrightarrow x<y\ \text{ou}}
\]
% archive sheet number: 91

